\section{Introduction}

Modern cloud environments provide hundreds of services on various platforms
that are readily available for developers to build a
large-scale distributed application. System architects and developers
following the microservice or service-oriented designs can freely
choose and compose the building blocks in a modular fashion and plug them
into the application.

However, different distributed storage systems provide different consistency
semantics, and the mismatch of the chosen system's semantics and the
application's expected semantics could introduce subtle and silent
semantic bugs \cite{silentbugs, silentbugs2}. For example, a user may connect a
key-value store providing
eventual consistency to an application that expects a
linearizable storage service. In this case, the application could
read outdated information causing erroneous outcomes but without
crash-like explicit errors.

Developers and system architects could carefully review the
application needs and choose a
matching system. However, when the application depends on multiple
distributed services, checking the overall compatibility between the
application logic and multiple system components can easily become unmanageable. Even
if one successfully finds the correct combination, for example, through
extensive testing, developers might change the application code or
plug in different systems due to management, considerations of cost and performance.
Consequently, sophisticated compatibility checks must be repeated.

% When you have fixed application logic, you may choose two
% compatible systems. Two systems may collectively satisfy the
% application requirement, and changing one may invalidate both.
% NOTE: addressed below

Formal verification offers an alternative approach to these challenges,
where the correctness of the code can mathematically guarantee the
compatibility between code and system \cite{lampson2021hints}.
However, depending on
the complexity of the code, which could include non-essential parts
for reasoning about compatibility, the mechanized formal proofs
can take extensive time and effort \cite{certikos-osdi16,
klein2009sel4, wormspace}.
Furthermore, depending on the technique used, mechanized formal
proofs can be several times larger, in Lines of Code (LoC), than the implementation.
Moreover, minor changes can easily invalidate the entire proof
\cite{klein2009sel4}.
Thus, formal verification
leads to a similar problem to the above-mentioned testing methods where
the compatibility check
lacks the agility to adapt to changes in underlying systems and
applications.

In this paper, we propose \sysname{}, a lightweight \sysnamelong{}. \sysname{}
helps developers quickly encode their application logic and
requirements and check the compatibility with underlying distributed
systems using an SMT-based approach.
%to encode and analyze
%the safety properties of different consistency semantics and
%their compositions.

\sysname{} includes axiomatic models of distributed consistency
semantics, including eventual consistency,
linearizability~\cite{viotti2016consistency}, four session
guarantees~\cite{terry1994session}, pipeline RAM (PRAM) consistency~\cite{pram},
and causal consistency~\cite{causaldef}, which are commonly used in both
industry and academia~\cite{pileus, cosmosdb, anna, bayou, yogurt,
pnuts, cloudflare-session, causalmemory, cops}.
Users of our tool can leverage predefined models as
building blocks to rapidly iterate the service compositions and compatibility
of these models with their applications. Additional models can be
added to further
strengthen \sysname{}'s expressiveness with provided interfaces.
User-authored custom models inheriting from our base class automatically work with our pre-built ones,
including assertions for individual satisfiability and compatibility between one or more models
with user-defined application logic.

Although axiomatic definitions of distributed consistency semantics
exist in prior literature \cite{finkbeiner2007smt,
  tekken2017monitoring, kuper2019toward, fan2023satisfiability,
ma2024noctua}, encoding the details into the SMT
clauses for compatibility checks is challenging as the definitions
omit key details. For example, the commonly used concept
of \textit{visibility} for defining weak distributed consistency
semantics is subtle, but to our knowledge, no literature provides
a concrete definition. \sysname{} fills in these missing gaps to realize
full-fledged formal models for the compatibility checking.

Each application functionality is systematically encoded as a
directed acyclic graph (DAG), and the entire application is modeled as
a directed multi-graph format which we call the \sysgraph{}.
Each node in the \sysgraph{} represents key functionality in the application
that requires certain properties from the lower-level nodes it depends on,
and provides certain properties to the upper-level nodes.
The edges are used to specify more specific requirements and conditions that each
functionality needs to enforce in accessing the downstream nodes.
The leaf nodes typically represent external distributed services
following one of the pre-built distributed consistency semantic models
in \sysname{}.
%
The \sysgraph{} model is expressive enough to model an application logic
from high-level and trace complex dependencies that go across multiple
application logic boundaries, and distributed systems which the application
relies on. 

Using the \sysgraph{}, \sysname{} can verify simple applications 
relying on one distributed service and scale to complex microservice-based
production environments that leverage multiple modules and distributed storage systems.
%
Once the user specifies the application logic, \sysname{} can quickly
check the compatibility with the plugged-in consistency models. This
could be verifying the user's selection, or alternatively, \sysname{}
can explore different combinations of existing models to
find compositions of distributed systems with compatible consistency semantics.
These approaches enable fast exploration and iteration of alternative
system designs, guiding the development to adhere to intended
consistency guarantees from the very beginning.

%%%%%% Examples %%%%%%    % NOTE: addressed here
%%%%%% Examples %%%%%%    {
%%%%%% Examples %%%%%%      \color{red}
%%%%%% Examples %%%%%%      Consider a Not So Simple Storage Service
% (NS4) that advertises
%%%%%% Examples %%%%%%      itself to provide linearizability
% guarantee to users' read/write to objects.
%%%%%% Examples %%%%%%      NS4's architecture is based on an API
% layer that provides
%%%%%% Examples %%%%%%      arbitration between users and the
%%%%%% Examples %%%%%%      underlying monolithic storage layer
%%%%%% Examples %%%%%%      with read-your-writes guarantee. In this
% case, \sysname{} can be used to
%%%%%% Examples %%%%%%      model the abstract semantical properties
% of NS4, and check whether the
%%%%%% Examples %%%%%%      modeled NS4 have advertised user-facing
% linearizable guarantee.
%%%%%% Examples %%%%%%      Three nodes (client, API, and storage)
% and two edges (client to
%%%%%% Examples %%%%%%      API, API to storage) are
%%%%%% Examples %%%%%%      added to the DAG, where the client node
% requires linearizability
%%%%%% Examples %%%%%%      from the API node,
%%%%%% Examples %%%%%%      API node exports linearizability to the
% client, and API node
%%%%%% Examples %%%%%%      requires read-your-writes from the storage node.
%%%%%% Examples %%%%%%      The storage node exports read-your-writes
% to the API node. With
%%%%%% Examples %%%%%%      this modeling, \sysname{}
%%%%%% Examples %%%%%%      will be able to find a satisfying
% assignment for the DAG where
%%%%%% Examples %%%%%%      user-facing guarantee is linearizability.
%%%%%% Examples %%%%%%      However, if we swap the API layer
% requirement from the storage
%%%%%% Examples %%%%%%      layer from read-your-writes to
%%%%%% Examples %%%%%%      to PRAM consistency, \sysname{} will be
% able to find that the leaf
%%%%%% Examples %%%%%%      level semantical requirement can no
% longer be satisfied.
%%%%%% Examples %%%%%%    }

%%%%%% %%%%%%    \jiyong{Some challenges in doing this?}
%%%%%% %%%%%%    \yifei{What aspect should we talk about? Preliminary
% encoding on
%%%%%% %%%%%%      semantics? Debugging? Unresolved encoding on
% models with "eventual"
%%%%%% %%%%%%    visibility (see eventual source code)?}

As case studies, we use \sysname{} to model two microservice-based
applications \cite{gan2019an}, a movie streaming service, an e-commerce
service, and a cross-service causal storage system built on top of
multiple weakly consistent distributed storage systems \cite{boltoncc}.
We evaluate and demonstrate \sysname{}'s ability to check compatibility
between user-defined application specific logic, with underlying
distributed systems and explore other compatible semantics.
Our framework is capable of checking a single compatibility combination in under 200 milliseconds and exploring
other compatible combinations in less than 5 minutes.

%%%%%% %%%%%%    \jiyong{Some insights we gained?}
%%%%%% %%%%%%    \yifei{
%%%%%% %%%%%%      This is a project with a lot of moving parts,
% modeling a system,
%%%%%% %%%%%%      like in real life, have multiple ways to achieve
% the same goal.
%%%%%% %%%%%%      The level of abstraction we've chosen made it not
% too hard to check
%%%%%% %%%%%%      real system's behavior but that part is not implemented.
%%%%%% %%%%%%      Personally I feel like some part of the project
% would be better
%%%%%% %%%%%%      written in a proper theorem prover (representation
% of ops, binary
%%%%%% %%%%%%      relations), and discharge part of semantic
% constraints to SMT solvers.
%%%%%% %%%%%%    }

This paper makes the following contributions:
\itemparen{
%\item We present \sysname{}, the first compatibility checking tool for
%  applications that are composed of multiple distributed services.
\item We present \sysname{}, the first semantic-compatibility-checking tool 
	between an application and the multiple distributed services that it relies on.
\item We provide the missing formal definitions including visibility to fully express
  and check distributed weak consistency semantics.
\item We present a systematic way of encoding complex applications
  using the \sysgraph{} structure and checking their compatibility
  collectively using the SMT-based approach.
\item We conduct case studies on three representative distributed
  applications to demonstrate the workings of \sysname{}.
}

%%%%    % FIXME: full rewrite
%%%%    {\color{red}
%%%%
%%%%      Modern distributed systems often relax strict consistency
%%%%      guarantees to achieve
%%%%      scalability, fault tolerance, and low latency operations.
% Understanding and
%%%%      comparing different forms of relaxed consistency semantics
% is crucial for both
%%%%      system designers and researchers. Yet, systematically
% modeling and reasoning
%%%%      about these semantics, as well as their compositional
% properties, remained a
%%%%      understudied area.
%%%%
%%%%      Our framework aims to bridge this gap by providing a
% SMT-based approach to
%%%%      encode and analyze the safety properties of weaker consistency
%%%%      models (comparing
%%%%      with linearizability \cite{herlihy1990linearizability}), and their
%%%%      compositions under
%%%%      the assumption of livenss properties are always satisfied,
% i.e. eventual
%%%%      consistency \cite{petersen1997flexible} is guaranteed (replicas
%%%%      eventually converge),
%%%%      and we are layering everything on top of it. By offering a per-system
%%%%      (differentiated by symbols), logical representation of
% operations, our system
%%%%      can differentiates operations belonging to distinct systems
% and/or components
%%%%      through the use of unique symbols (see
% tests/test\_cross\_causal.py), enabling
%%%%      analysis of cross-system interactions. Additionally, solver
% backends can
%%%%      discharge proof witnesses on the satisfiability of specific
% consistency
%%%%      properties, or the compatibility and composability between
% multiple models.
%%%%      These approaches enable fast exploration and iteration of
% alternative system
%%%%      designs, guiding the development adhere to intended consistency
%%%%      guarantees from
%%%%      the very beginning.
%%%%
%%%%    }
